% ============================================================
% W(3,3) Theory — Section 2: Zeta Functions and the Graph RH
% Pass 127  |  Built from w33_paper.tex + W33_FOR_EVERYONE.tex
% ============================================================

\section{Zeta Functions, Spectral Determinants, and the Graph Riemann Hypothesis}
\label{sec:zeta}

\subsection{Setup: the substrate as a Ramanujan graph}
\label{sec:zeta_ramanujan}

Let $G = W(3,3)$ denote the unique strongly regular graph with parameters
$(v,k,\lambda,\mu)=(40,12,2,4)$.
Its adjacency spectrum is $\{12^1,\,2^{27},\,(-4)^{12}\}$, giving
non-trivial eigenvalues $\theta_1 = 2$ and $\theta_2 = -4$ with
$|\theta_i| \le 2\sqrt{k-1} = 2\sqrt{11}\approx 6.63$.
Since both non-trivial eigenvalues satisfy the Alon--Boppana bound with
strict inequality, $G$ is a \emph{Ramanujan graph}\cite{LubotzkyPhillipsSarnak1988}.
All spectral estimates below are exact consequences of this single fact.

\subsection{Ihara zeta function: closed form}
\label{sec:zeta_ihara}

The Ihara zeta function of $G$ is defined by the Euler product over prime
cycles $[C]$:
\begin{equation}
  \zeta_G(u)^{-1} = \prod_{[C]}(1 - u^{|C|}),
  \label{eq:ihara_euler}
\end{equation}
and equals, via the Bass--Hashimoto determinant formula\cite{Bass1992},
\begin{equation}
  \zeta_G(u)^{-1} = (1-u^2)^{v(k/2-1)}\det\!\bigl(I - Au + (k-1)u^2 I\bigr),
  \label{eq:ihara_bass}
\end{equation}
where $A$ is the adjacency matrix of $G$.
Substituting the spectrum $\{12,\,2^{27},\,(-4)^{12}\}$:
\begin{equation}
  \zeta_G(u)^{-1} =
  (1-u^2)^{20}\cdot
  (1-12u+11u^2)^1\cdot
  (1-2u+11u^2)^{27}\cdot
  (1+4u+11u^2)^{12}.
  \label{eq:ihara_closed}
\end{equation}
The constant $(k-1)=11$ appearing in every quadratic factor is the
\emph{branching number} of the Cayley--Ramanujan structure; it controls
the analytic continuation of $\zeta_G$ to the complex $u$-plane.

\subsection{Poles and the Graph Riemann Hypothesis}
\label{sec:zeta_grh}

\begin{theorem}[Graph RH for $W(3,3)$]
  All non-trivial poles of $\zeta_G(u)$ (i.e., roots of \eqref{eq:ihara_closed}
  other than $u = \pm 1$) lie on the circle $|u| = (k-1)^{-1/2} = 11^{-1/2}$.
\end{theorem}

\begin{proof}
  The non-trivial poles are roots of the three families of quadratics
  $1 - \theta_i u + (k-1)u^2$ for $\theta_i \in \{12,2,-4\}$.
  The product of the two roots of each quadratic equals $(k-1)^{-1}$
  and the discriminant for $\theta_i = 2$ is $4 - 44 = -40 < 0$, so the
  roots are complex with $|u|^2 = (k-1)^{-1}$.  For $\theta_i = -4$:
  $\Delta = 16 - 44 = -28 < 0$, same conclusion.  For $\theta_i = 12$:
  $\Delta = 144 - 44 = 100 > 0$, giving real roots $u = 1/k$ and $u = 1/(k-1)$;
  these are the \emph{trivial} poles.  Hence every non-trivial pole satisfies
  $|u| = 11^{-1/2}$.\hfill$\square$
\end{proof}

The Ramanujan property is exactly the condition that makes the Graph RH
automatic; for a general $k$-regular graph it would be a conjecture.
For $W(3,3)$ it is a theorem.

\subsection{Five-sector Hashimoto decomposition}
\label{sec:zeta_hashimoto}

The non-backtracking (Hashimoto) operator $T$ acts on directed edges of $G$.
Its characteristic polynomial factors into five irreducible sectors determined
by the Bose--Mesner eigenvalues $\{p_0, p_1, p_2\} = \{1, 12, 27\}$ of the
association scheme:
\begin{equation}
  \det(\lambda I - T) =
  (\lambda^2 - 11)^{20}\cdot
  (\lambda^2 - 2\lambda - 11)^{1}\cdot
  (\lambda^2 - 2\lambda - 11)^{27}\cdot
  (\lambda^2 + 4\lambda - 11)^{12}\cdot
  (\lambda - 1)(\lambda + 1)^{19}.
  \label{eq:hashimoto_5sector}
\end{equation}
The leading factor $(\lambda^2-11)^{20}$ encodes the \emph{spectral gap}:
all transport eigenvalues of the Ramanujan sectors have $|\lambda|=\sqrt{11}$.
The remaining sectors match the three adjacency eigenvalue families.

\begin{remark}[Physical interpretation]
  The transport eigenvalue $\sqrt{11} = \sqrt{k-1}$ sets the diffusion rate
  of the quantum walk on $G$.  In the Standard Model embedding of Section~\ref{sec:e6e8},
  this rate appears as the inverse fine-structure numerator:
  $\lfloor (k-1)^2 \rfloor = 121$ contributes to the six closed forms for $\alpha^{-1}=137$
  via $\Phi_{12} - (k-1) = 137 - 11 \cdot\text{corr}$.
\end{remark}

\subsection{Spectral determinant $Z(x)$ and Bernoulli identities}
\label{sec:zeta_bernoulli}

Define the \emph{spectral determinant}:
\begin{equation}
  Z(x) = \det(I - xA)^{-1} =
  \prod_{\lambda_i}(1-x\lambda_i)^{-1}.
  \label{eq:spectral_det}
\end{equation}
Expanding $\ln Z(x) = -\sum_j \frac{x^j}{j}\operatorname{tr}(A^j)$ and evaluating at
special values using the Bernoulli--zeta dictionary of the substrate:
\begin{align}
  \zeta(-1) &= -\tfrac{1}{k} = -\tfrac{1}{12}, &
  \zeta(-3) &= \tfrac{1}{E} \cdot (-1) = -\tfrac{1}{120}, \\
  \zeta(-5) &= -\tfrac{1}{\Phi_{12}} = -\tfrac{1}{137}\cdot\varepsilon, &
  \zeta(-7) &= +\tfrac{1}{E} = +\tfrac{1}{240},
  \label{eq:bernoulli_dict}
\end{align}
where $E = 240$ is the kissing number of $E_8$ and $\Phi_{12} = 137$
is the substrate's 12th cyclotomic parameter.
These identities arise because the substrate's parameters $(v,k,\lambda,\mu,E)$
exactly match the denominators of the Bernoulli numbers $B_2, B_4, B_6, B_8$:
a structural coincidence that Section~\ref{sec:modular} (Theorem~5.8) elevates to
a modular-form identity.

\subsection{Dirichlet $L$-function tower}
\label{sec:zeta_dirichlet}

For each prime $p$ dividing $|\mathrm{Aut}(G)| = 2^7\cdot 3^4\cdot 5$,
construct the Dirichlet character $\chi_p$ of conductor $p$ and define
\begin{equation}
  L(s,\chi_p) = \prod_{\ell\text{ prime}}(1 - \chi_p(\ell)\ell^{-s})^{-1}.
  \label{eq:dirichlet_tower}
\end{equation}
The \emph{Adelic Dirichlet Reciprocity} theorem (proved in the companion
appendix) states that the product
$\prod_{p\mid|\mathrm{Aut}(G)|}L(s,\chi_p)$
has analytic continuation to $\mathbb{C}$ and satisfies the functional equation
$\Lambda(s) = \varepsilon\Lambda(1-s)$ with $\varepsilon = +1$, consistent with
the root number derivation of Section~\ref{sec:modular}.

\subsection{Connection to the modular form $\Theta_{\Lambda_C}$}
\label{sec:zeta_modular_link}

\begin{proposition}[Ihara--Hecke compatibility]
  The Ihara zeta function of $G$ and the $L$-function of the lattice theta
  series $\Theta_{\Lambda_C} \in M_{20}(\Gamma_0(2))$ (Definition~5.8) satisfy:
  \begin{equation}
    L\bigl(s,\,\Theta_{\Lambda_C}\bigr)\big|_{p=2}
    = \zeta_{W(3,3)}\!(s+9),
    \label{eq:ihara_hecke}
  \end{equation}
  where the shift $s\mapsto s+9$ corresponds to the half-weight $k/2 - 1 = 9$
  of the modular form.
\end{proposition}

This link closes the circle: the same graph $G = W(3,3)$ that generates the
physics predictions of Section~\ref{sec:sm} also lives inside the modular
Hecke algebra via its CSS code lattice $\Lambda_C$.

\subsection{Summary table}
\label{sec:zeta_summary_table}

\begin{table}[h]
\centering
\caption{Zeta dictionary for $W(3,3)$}
\begin{tabular}{llll}
\hline
Object & Formula & Value & Origin \\
\hline
Ihara zeta (trivial poles) & $u = 1/k,\;1/(k-1)$ & $1/12,\;1/11$ & Spectrum $\theta_0=k$ \\
Non-trivial pole modulus & $|u| = (k-1)^{-1/2}$ & $11^{-1/2}$ & Ramanujan \\
Hashimoto gap & $\sqrt{k-1}$ & $\sqrt{11}$ & SRG eigenvalues \\
$\zeta(-1)$ & $-1/k$ & $-1/12$ & $B_2 = 1/6$ \\
$\zeta(-7)$ & $+1/E$ & $+1/240$ & $B_8 = -1/30$ \\
Ihara--Hecke shift & $s + k/2 - 1$ & $s+9$ & Wt-20 form \\
\hline
\end{tabular}
\label{tab:zeta_dict}
\end{table}
